\begin{frame}[fragile]{Multi-Thread Language}

\begin{itemize}
	\item Java is a multi-thread programming language.
	\item We can extend the Thread class (lines 1 -- 6) to instantiate them (line 12) and execute the run() method in a thread (line 13). 
\end{itemize}

\begin{Code}{}%
\begin{lstlisting}[firstnumber=1, label=glabels, xleftmargin=0pt, basicstyle=\tiny]% 

class Multi extends Thread {
  public void run()
  {
    try { ... }
    catch (Exception e) { ... }
  }
}
public class Multithread {
  public static void main(String[] args) {
    int n = 8; // Number of threads
    for (int i = 0; i < n; i++) {
      Multi object = new Multi();
      object.start();
    }
  }
}
\end{lstlisting}%   
\end{Code}%
\end{frame}

\begin{frame}[fragile]{}

\begin{itemize}
	\item Most languages support this multi-thread programming model. 
	\item These are examples from C\#, but other programming languages such as C++ or Rust are almost identical. 
\end{itemize}

\begin{Code}{}%
\begin{lstlisting}[firstnumber=1, label=glabels, xleftmargin=0pt, basicstyle=\scriptsize]% 

using System;
using System.Threading;
 
namespace Threads
{
  class Program {
    public static void Main(string[] args) {
      Thread thread = new Thread(new ThreadStart(Task));
      thread.Start();
      ...
      thread.Join();
   }
   private static void Task() { ... }
}   
\end{lstlisting}%   
\end{Code}%
\end{frame}

